% =====================================================================
% 第 III 章  Class invariants I (A. Borel, Nov. 6, 1957)
% =====================================================================
\lecture{Class invariants I}{Class Invariants I}{A. Borel}{Nov. 6, 1957}

\setcounter{theorem}{0}
\renewcommand{\thetheorem}{\arabic{theorem}}

% ---------------------------------------------------------------
\section{Introduction}\label{sec:III-intro}

Let $X=\mathbb{C}/\Gamma$ be an elliptic curve. If $X$ has a non-trivial
complex multiplication, the ring $A(X)$ of all complex multiplications of
$X$ forms an order of an imaginary quadratic field. Let an imaginary
quadratic field $K$ be fixed. Then there is a one-one correspondence
between the set of ideal classes $\mathfrak{k}_{1},\ldots,\mathfrak{k}_{n}$
of $K$ and the set of $X$ such that $A(X)$ coincides with the ring of all
algebraic integers in $K$; the correspondence is given by assigning to
$\mathfrak{k}_{i}$ the elliptic curve with the invariant
$j(\mathfrak{k}_{i})$. (Cf.\ chapter I, §1.)

The main purpose of the four following lectures is to prove that the class
invariants $j(\mathfrak{k}_{i})$ are all algebraic integers forming a full
set of conjugates over $K$, and also over $\mathbb{Q}$, and that
$K(j(\mathfrak{k}_{i}))$ is the maximal unramified abelian extension of
$K$. As mentioned in chapter I, we follow Hasse's first method, which uses
analysis and class field theory.

% ---------------------------------------------------------------
\section{Modular correspondences}\label{sec:III-corr}

Let $E$ be the upper-half complex plane, $G$ the modular group and $S$ the
closed Riemann surface obtained by adjoining $\infty$ to $E/G$. $S$ is
mapped isomorphically onto the complex sphere by the function $j$ (cf.\ II,
§3). So $S$ is an algebraic curve of genus $0$ and an algebraic
correspondence of $S$ into itself is called a \emph{modular
correspondence}.

A method to generate such a modular correspondence is as follows:

Let $\mathcal{G}$ be the group of all automorphisms of the Riemann surface
$E$. It is the quotient of the group $\mathrm{GL}(2,\mathbb{R})$ of real
$2\times 2$ matrices by its center. Take a subset $H$ of $\mathcal{G}$
containing $G$ such that $GH=H$, $[H:G]=m<\infty$ and $H^{-1}=H$ (hence
also $HG=H$). For any $M$ in $H$, let $j_{M}(z)$ denote the function
$j(M(z))$. Then $j_{M}$ depends only upon the class $GM$ of $M$ in $H/G$.
Let $j_{1}(z),\ldots,j_{m}(z)$ be the functions obtained in this way from
all $m$ classes of $H/G$. Then it can easily be proved that each element
of $H$ may be represented by a matrix with integral rational coefficients
and that each class $GM$ contains a triangular matrix. It follows then
(see the proof of Theorem~1a below) that the symmetric functions of
$j_{1},\ldots,j_{m}$ are integral modular functions, i.e.\ polynomials in
$j$. Hence, $F(t,j)=\prod_{s=1}^{m}(t-j_{s}(z))$ is a polynomial of $t$
over $\mathbb{C}(j)$, and a mapping which maps a point of $S$ with
``coordinate'' $j(z)$ to the points with coordinates
$j_{1}(z),\ldots,j_{m}(z)$ defines a modular correspondence.

Taking any modular function $g(z)$ instead of $j(z)$, we can define
$g_{1}(z),\ldots,g_{m}(z)$ just as above, and symmetric functions of
$g_{1},\ldots,g_{m}$ then give us modular functions. More generally, let
$h$ be any modular form of weight $k$. For $M$ in $H$, put
$h_{M}(z)=h(M(z))\bigl(\frac{dM(z)}{dz}\bigr)^{k}$. Then $h_{M}(z)$ again
depends only upon the class $GM$ and we get $m$ such functions
$h_{1}(z),\ldots,h_{m}(z)$ from $m$ classes of $H/G$. A symmetric function
of weight $\nu$ in $h_{1},\ldots,h_{m}$ is then a modular form of weight
$\nu k$.

% ---------------------------------------------------------------
\section{\texorpdfstring{The correspondences $F_{n}$}{The correspondences Fn}}\label{sec:III-Fn}

We now consider a special kind of correspondences. Let $n$ be any positive
integer. We denote by $H_{n}^{*}$ the set of all automorphisms of $E$
given by
\[
M:\ z\longrightarrow\frac{az+b}{cz+d},\qquad
ad-bc=n,\quad a,b,c,d\in\mathbb{Z},
\]
and denote by $H_{n}$ the subset of $H_{n}^{*}$ consisting of those $M$
satisfying $(a,b,c,d)=1$. Clearly, $GH_{n}^{*}=H_{n}^{*}$,
$(H_{n}^{*})^{-1}=H_{n}^{*}$ and $GH_{n}=H_{n}$, $(H_{n})^{-1}=H_{n}$.
Furthermore, a set of representatives for classes $GM$ of $H_{n}^{*}/G$ is
obtained by the automorphisms with matrices
\[
\begin{pmatrix}a&b\\ 0&d\end{pmatrix},\qquad
a>0,\ d>0,\ d>b\geq 0,\ ad=n.
\tag{1}\label{eq:III-1}
\]
Hence, $[H_{n}^{*}:G]=\sigma_{1}(n)=\sum_{d\mid n}d$. Imposing a further
condition $(a,b,d)=1$ on matrices in \eqref{eq:III-1}, we get a set of
representatives for $H_{n}/G$ and it is not difficult to see that
$[H_{n}:G]=\psi(n)=n\prod_{p\mid n}\bigl(1+\frac{1}{p}\bigr)$.

Another way of defining $H_{n}$ is as follows: Let $\Gamma'$ be a
two-dimensional lattice on the complex plane. A sublattice $\Gamma'$ of
$\Gamma$ with index $n$ is called primitive if $\Gamma/\Gamma'$ is a
cyclic group of order $n$. Let $\omega_{1},\omega_{2}$ be a basis of
$\Gamma$ and $\omega_{1}',\omega_{2}'$ a basis of such a primitive
sublattice $\Gamma'$. Put
$\omega_{1}'=a\omega_{1}+b\omega_{2}$,
$\omega_{2}'=c\omega_{1}+d\omega_{2}$. Then
$M:z\longrightarrow\frac{az+b}{cz+d}$ is in $H_{n}$. Conversely, every $M$
in $H_{n}$ can be obtained in this way from some $\Gamma,\Gamma'$ and
their suitable bases $\omega_{1},\omega_{2}$ and
$\omega_{1}',\omega_{2}'$.

Thus if $X$ is a curve with lattice of periods $\Gamma$, the
correspondence $F_{n}$ defined by $H_{n}$ associates to $X$ the curves
having a lattice of periods primitive of index $n$ in $\Gamma$.
Analogously, $H_{n}^{*}$ associates to $X$ all curves which have a period
lattice of index $n$ in $\Gamma$.

Given such $\Gamma'$ primitive of index $n$ in $\Gamma$, we can always
find, by a theorem on abelian groups, a basis $\omega_{1},\omega_{2}$ of
$\Gamma$ such that $n\omega_{1}$, $\omega_{2}$ form a basis of $\Gamma'$.
Hence, if $M_{0}$ denotes the automorphism $z\longrightarrow nz$, then
$H_{n}=GM_{0}G$. Therefore, right multiplication of $G$ on $H$ permutes
the classes $GM$ of $H/G$ transitively. In general, this is not true for
$H_{n}^{*}$.

% ---------------------------------------------------------------
\section{The modular equations}\label{sec:III-modeq}

Now, let $M_{1},\ldots,M_{N}$ $(N=\psi(n))$ be the representatives for
$H_{n}/G$ given by the matrices \eqref{eq:III-1} with $(a,b,d)=1$. Put
\[
j_{s}(z)=j\bigl(M_{s}(z)\bigr),\qquad s=1,\dots,N.
\]
These are the functions associated with $N$ classes of $H_{n}/G$ and, as
$H_{n}G=H_{n}$, $j_{s}(z)\longrightarrow j_{s}(T(z))$ is a permutation of
$j_{1},\ldots,j_{N}$ for any $T$ in $G$. Hence, a symmetric function of
$j_{1},\ldots,j_{N}$ is invariant under $z\longrightarrow T(z)$,
$T\in G$.

\begin{theorem}\label{thm:III-1}
(a) $F_{n}(t,j)=\prod_{s=1}^{N}(t-j_{s}(z))$ is a polynomial of $t$ and
$j=j(z)$ with integral rational coefficients;

(b) if $n$ is not a square, the highest coefficient of $j$ in
$F_{n}(j,j)$ is $\pm 1$;

(c) $F_{n}(t,j)$ is an irreducible polynomial of $t$ over the field
$\mathbb{C}(j)$;

(d) $F_{n}(t,j)=F_{n}(j,t)$, $n>1$.
\end{theorem}

\begin{proof}[Proof.]
(a) An elementary symmetric function
$\sigma_{\nu}(j_{1}(z),\ldots,j_{N}(z))$ of $j_{1},\ldots,j_{N}$ is, as
noticed above, invariant under $G$ and is obviously a holomorphic function
of $z$ in $E$. To see the behavior of $\sigma_{\nu}$ at $\infty$, put
$q=e^{2\pi iz}$. Then, by II §6,
$j(z)=q^{-1}(1+A(q))$ where $A(q)$ is a power series of $q$ with integral
rational coefficients and $A(0)=0$. Let
\[
M_{s}\leftrightarrow
\begin{pmatrix}a_{s}&b_{s}\\ 0&d_{s}\end{pmatrix},\qquad
e^{2\pi iM_{s}(z)}=\zeta_{d_{s}}^{b_{s}}q^{\frac{a_{s}}{d_{s}}},\qquad
\zeta_{d_{s}}=e^{\frac{2\pi i}{d_{s}}}.
\]
Then,
\[
j_{s}(z)=j\bigl(M_{s}(z)\bigr)
=\zeta_{d_{s}}^{-b_{s}}\,q^{-\frac{a_{s}}{d_{s}}}
\Bigl(1+A\Bigl(\zeta_{d_{s}}^{b_{s}}q^{\frac{a_{s}}{d_{s}}}\Bigr)\Bigr).
\tag{2}\label{eq:III-2}
\]
Therefore, $\sigma_{\nu}$ can only have the singularity of a pole in $q$
at $q=0$, and, hence, is a polynomial in $j$ over $\mathbb{C}$.
Furthermore, the coefficients of the $q$-expansion of $\sigma_{\nu}$ are
all algebraic integers in the field $\mathbb{Q}(\zeta_{n})$,
$\zeta_{n}=e^{2\pi i/n}$. Let $\tau$ be any Galois automorphism of
$\mathbb{Q}(\zeta_{n})/\mathbb{Q}$. Then $\tau(\zeta_{n})=\zeta_{n}^{f}$
for some $f$ with $(n,f)=1$, and we can write
\[
\tau\bigl(\zeta_{d_{s}}^{b_{s}}\bigr)=\zeta_{d_{t}}^{fb_{s}}
=\zeta_{d_{t}}^{b_{t}},
\]
where
\[
\begin{pmatrix}a_{t}&b_{t}\\ 0&d_{t}\end{pmatrix},\qquad
a_{t}=a_{s},\quad d_{t}=d_{s},\quad b_{t}\equiv fb_{s}\bmod d_{s},
\]
is another matrix in \eqref{eq:III-1} uniquely determined by $M_{s}$.
Hence, $\tau(j_{s})=j_{t}$ and $\tau$ permutes the functions
$j_{1},\ldots,j_{N}$ among themselves. Thus,
$\tau(\sigma_{\nu})=\sigma_{\nu}$ for any $\tau$ in the Galois group of
$\mathbb{Q}(\zeta_{n})/\mathbb{Q}$, and the coefficients of the
$q$-expansion of $\sigma_{\nu}$ must be rational integers. By the
q-expansion principle (II §6), $\sigma_{\nu}$ is then a polynomial in $j$
with integral rational coefficients.

(b) Since $n$ is not a square, $\frac{a_{s}}{d_{s}}$ in
\eqref{eq:III-2} cannot be $1$. Hence, the leading coefficient of the
$q$-expansion of $j-j_{s}$ is a root of unity, and so is the leading
coefficient of the q-expansion of the product
$F_{n}(j,j)=\prod_{s=1}^{N}(j-j_{s})$. However, this coefficient is equal
to the highest coefficient of $j$ in $F_{n}(j,j)$ and, as it is rational,
it must be $+1$.

(c) This follows immediately from the fact that $j_{1},\ldots,j_{N}$ are
permuted transitively among themselves under
$j_{s}(z)\longrightarrow j_{s}(T(z))$, $T\in G$. (Cf.\ the last lines of
§\ref{sec:III-Fn}.)

(d) Since $z\longrightarrow nz$ and $z\longrightarrow z/n$ are both
contained in $H_{n}$,
$F_{n}(j(nz),j(z))=0$ and $F_{n}(j(z/n),j(z))=0$, identically in
$z\in E$. Replacing $z$ by $nz$ in the second equality, we get
$F_{n}(j(z),j(nz))=0$. Hence, as polynomials in $\mathbb{C}(j)[t]$,
$F_{n}(t,j)$ and $F_{n}(j,t)$ have a common zero $t=j(nz)$. Since
$F_{n}(t,j)$ is irreducible by (c) and since the highest coefficient of
$t$ in $F_{n}(t,j)$ is $1$, we obtain
\[
F_{n}(j,t)=P(t,j)\,F_{n}(t,j),
\]
with a polynomial $P(t,j)$ in $\mathbb{Z}[t,j]$. It then follows that
\[
F_{n}(t,j)=P(j,t)\,F_{n}(j,t)=P(j,t)\,P(t,j)\,F_{n}(t,j),
\]
and hence, that
\[
P(t,j)=P(j,t)=+1.
\]
If $P=-1$, we would have $F_{n}(j,t)=-F_{n}(t,j)$ and $F_{n}(j,j)=0$, so
that $F_{n}(t,j)$ is divisible by $t-j$, contradicting (c). Hence, $P=1$
and $F_{n}(t,j)=F_{n}(j,t)$.
\end{proof}

$F_{n}(t,j)=0$ is the modular equation for the degree $n$. (For a given
$j(z)$, which is the invariant of an elliptic curve $X$, the roots of
$F_{n}(t,j)=0$ are the modular invariants of the curves $X_{i}$
($1\leq i\leq N$) which the correspondence $F_{n}$ associates to $X$.
Theorem~\ref{thm:III-1}d means that then $X$ is among the images of the
$X_{i}$; this can also be seen geometrically by noticing that if $\Gamma'$
is primitive of index $n$ in $\Gamma$, then $n\Gamma$ is primitive of
index $n$ in $\Gamma'$, and that the curves with periods $\Gamma$ and
$n\Gamma$ are isomorphic.)

% ---------------------------------------------------------------
\section{Class invariants}\label{sec:III-classinv}

We next show that an elliptic curve $X$ has non-trivial complex
multiplications if and only if the invariant $j(X)$ satisfies
$F_{n}(j(X),j(X))=0$ for some $n>1$. Let $X=\mathbb{C}/\Gamma$ and let
$\omega_{1},\omega_{2}$ be a basis of $\Gamma$ so that $j(X)=j(\omega)$,
$\omega=\omega_{1}/\omega_{2}$. Suppose, first, that
$F_{n}(j(\omega),j(\omega))=0$, $n>1$. Then $j_{s}(\omega)=j(\omega)$ for
some $s$ and $M(\omega)=\omega$ for some $M$ in $H_{n}$. Hence, there
exists a $w$ in $\mathbb{C}$ such that
\[
\begin{aligned}
w\omega_{1}&=a\omega_{1}+b\omega_{2},\\
w\omega_{2}&=c\omega_{1}+d\omega_{2},
\end{aligned}
\qquad
M\leftrightarrow
\begin{pmatrix}a&b\\ c&d\end{pmatrix}.
\]
Then, $w$ cannot be a rational integer and $X$ has a non-trivial complex
multiplication. Conversely, suppose $X$ has non-trivial complex
multiplications. Then the ring of complex multiplications $A(X)$ is an
order of an imaginary quadratic field $K$ and we can find a $w$ in $A(X)$
such that $N_{K/\mathbb{Q}}(w)=n>1$ and that $w$ is not in any
$mA(X)$, $m>1$. Reversing the above argument, we see that $j(X)$ satisfies
$F_{n}(j(X),j(X))=0$. (Cf.\ below.)

We now fix an imaginary quadratic field $K$ and consider the class
invariants $j(\mathfrak{k}_{i})$, $i=1,\ldots,h$, of $K$, i.e.\ the
invariants $j(X)$ of elliptic curves $X$ of which $A(X)$ coincides with
the ring of all algebraic integers in $K$. Let $w$ be an algebraic integer
in $K$ such that $n=N_{K/\mathbb{Q}}(w)$ is a square-free integer $>1$.
Such a $w$ always exists; if $K=\mathbb{Q}(\sqrt{-1})$, take
$w=1+\sqrt{-1}$, and if $K=\mathbb{Q}(\sqrt{-m})$, $m>1$ square-free, take
$w=\sqrt{-m}$. Let $\omega_{1},\omega_{2}$ be a basis of an ideal in an
ideal class $\mathfrak{k}$ of $K$ and put
\[
\begin{aligned}
w\omega_{1}&=a\omega_{1}+b\omega_{2},\\
w\omega_{2}&=c\omega_{1}+d\omega_{2},
\end{aligned}
\qquad a,b,c,d\in\mathbb{Z}.
\]

Then $ad-bc=N_{K/\mathbb{Q}}(w)=n$ and, as $n$ is square-free,
$(a,b,c,d)=1$. Hence,
$M:z\rightarrow\frac{az+b}{cz+d}$ is in $H_{n}$. Since
$j(M(\omega))=j(\omega)$, the invariant $j(\mathfrak{k})=j(\omega)$ then
satisfies $F_{n}(j(\mathfrak{k}),j(\mathfrak{k}))=0$, and from the
previous theorem, we obtain immediately the following:

\begin{theorem}\label{thm:III-2}
The class invariants $j(\mathfrak{k}_{i})$, $i=1,\ldots,h$, of an
imaginary quadratic field are algebraic integers.
\end{theorem}
